\section{Positive interpolation and the strip exponent}\label{c6:strip-asymptotics}

Positive interpolation turns the finite strip functional into averages under a determinant density. We estimate its characteristic products and interpolation residuals to obtain $B_N\asymp N^{-1/4}$. 

\subsection{A positive determinant formula}
Retain the measure and spaces in \eqref{eq:legacy-P-1}, and put
\[
 \begin{gathered}
 c_1=C_{L/2},\quad c_3=C_{3L/2},\quad h_0=1/D,\quad h_1=c_1/D,\\
 F_0=M+\mathbf R c_1,\quad F=F_0+\mathbf R c_3,\quad G=M+\mathbf R g
 \end{gathered}
\]
(using real spans in the following). Here $N$ remains the strip width, while $d=\dim F_0=2N+1$ is the number of interpolation nodes. A complete jet at a label of multiplicity $k$ consists of derivatives of orders $0,\ldots,k-1$. Use jets of $C_{\sqrt\beta}(x)$ in $\beta$ as bases, ordered first by $\beta$, then by derivative order. Thus $M$ takes labels $0,L^2$ each with multiplicity $N$; the $g$ basis row in $G$ is inserted between. For $X=(x_1<\dots<x_d)$ in $(0,\infty)$, let $\nu$ denote the probability proportional to
\[
 \det F_0(X)\det G(X)\prod_i d\mu(x_i)
\]
(evaluation matrix with functions as rows). Indeed $g=\int_0^L C_t\,dt$, so the $G$ determinant integrates that for $G_t=M+\mathbf R C_t$, $0<t<L$. \begin{lemma}[Positive jet interpolation]\label{lemma:jet-interpolation}
All ordered jet determinants formed from complete jets of $C_{\sqrt\beta}(x)$ at nonnegative ordered labels are positive at distinct positive ordered points and remain nonnegative at confluent nonnegative points. Complete jet spaces are Chebyshev spaces on $[0,\infty)$, counting multiplicities.\end{lemma}
\begin{proof}
All ordered jet determinants are positive at ordered nonnegative points (also confluently, with derivatives at repeated points). This follows from the nonnegative Taylor minor expansion using $C_{\sqrt\beta}(x)=\sum_j \beta^j x^j/(2j)!$, with strictly positive first consecutive-degree term. Here generalized power determinants of increasing nonnegative integer powers at positive ordered points are strictly positive: a combination of $n$ powers has at most $n-1$ zeros counting multiplicity (divide by the first and differentiate successively/Rolle), and the sign is positive e.g. by coalescing near 1 (Wronskian). Confluent cases at zero are nonnegative by limits. In particular the spaces of complete jets are Chebyshev counting multiplicities on $[0,\infty)$. All Gram entries pairing a row from $F$ or $h_0,h_1$ and one from $G$ are finite (the density in $p=\sqrt{x}$ decays at rate $b=3L$).
\end{proof}

Write $e_y H=H(y)-I_{F_0,X}H(y)$ for evaluation minus interpolation at $X$ using $F_0$. Also use $E_y H=H(y)-I_{F,(0,X)}H(y)$. The determinant integrations below are instances of Andr\'eief's identity; see \cite[Equation~(1.3)]{TracyWidom1998}. We retain the interpolation calculation to fix its oblique projection and evaluation columns. Gram/Cramer identities using \eqref{eq:legacy-P-1} give, for $s_*=\rho/(4a_*)$ there,
\[
 \mathcal L(H)/(4a_*)=\nu[e_{-4}H]+s_*\nu[e_0 H].
\]
Indeed integrating a determinant product gives the determinant of the pairings by multilinearity, so the mean interpolant is the unique linear projection onto $F_0$ preserving the moments against $G$. Also
\[
 \nu[e_0 c_3]\,\nu[e_{-4}h_0]-\nu[e_{-4}c_3]\,\nu[e_0 h_0]
 =\nu[(e_0 c_3) E_{-4}h_0].
\]
For this append rows $c_3,h_0$ and evaluation columns $0,-4$ to the $F_0$ determinant and integrate against $\det G(X)$ (Schur complement). Since $d$ is odd, $e_0 c_3<0$ by ordered jet positivity. Let $\nu_3$ be the tilt of $\nu$ by $-e_0 c_3$. It will suffice to show
\begin{equation}
 \begin{gathered}
 \nu[-e_0 c_3]\gtrsim N^{-1/2},\qquad |\nu[e_0 h_i]|\le N^{-1+o(1)},\qquad
 |\nu[e_{-4}h_i]|=O(1)\quad (i=0,1),\\
 \nu_3[E_{-4} h_0]\asymp N^{-1/4}.
 \end{gathered}                                                     \label{eq:legacy-P-7}
\end{equation}
For the elimination, abbreviate
\[
 A_y=\nu[e_yc_3],\qquad H_{j,y}=\nu[e_yh_j]
 \quad(y=0,-4;\ j=0,1).
\]
The condition $\mathcal L(c_3-h_1)=0$ gives
\[
 s_*=-\frac{A_{-4}-H_{1,-4}}{A_0-H_{1,0}},
\]
whenever the denominator is nonzero. The determinant identity above, together with the definition of $\nu_3$, then gives the exact formula
\[
 \frac{\mathcal L(h_0)}{4a_*}
 =\frac{A_0\nu_3[E_{-4}h_0]
        +H_{1,-4}H_{0,0}-H_{1,0}H_{0,-4}}
       {A_0-H_{1,0}}.
\]
The estimates in \eqref{eq:legacy-P-7} imply that $|A_0|\gtrsim N^{-1/2}$ and $|H_{1,0}|\le N^{-1+o(1)}$. Hence the denominator is nonzero for large $N$, the coefficient of the positive residual is $1+O(N^{-1/2+o(1)})$, and the remaining quotient is $O(N^{-1/2+o(1)})$. The estimate for $\nu_3[E_{-4}h_0]$ therefore proves $B_N\asymp N^{-1/4}$.

\subsection{Characteristic products}

\begin{lemma}[Characteristic products]\label{lemma:characteristic-products}
The following characteristic-product estimates hold for both $\eta=\nu,\nu_3$. Put $r(x)=(x+y^2)/(x+z^2)$. We claim
\begin{equation}
 \eta\left[\prod_i r(x_i)\right]\asymp N^{b(y-z)/\pi}\qquad (0<y<z\text{ fixed}).                 \label{eq:legacy-P-8}
\end{equation}
In the separate growing-parameter regime $y^2=2z^2=2N^{3/2}$, where $1\le r\le2$, the expectation is at most $\exp(o(N))$. \end{lemma}

We isolate the elementary identity used in its proof, so that the polynomial comparison has an explicit starting point.
\begin{lemma}[Polynomial generating identity]\label{lemma:polynomial-generating}
Let $c>0$ and let $\mu_c$ be the probability measure proportional to $dp/\cosh(cp)$ on $p=\sqrt{x}>0$, and put $\alpha=\pi/(2c)$. Then
\begin{equation}
 A_m(x)=\cos(\alpha m)C_m(x),\quad w_m=\tan^2(\alpha m),\quad
 \int A_mA_n\,d\mu_c=1/(1-w_mw_n)
 \label{eq:legacy-P-9}
\end{equation}
for real $|m|+|n|<c$.
\end{lemma}
\begin{proof}
The elementary transform of $1/\cosh(cp)$ follows by integrating on the full line and shifting by $i\pi/c$, picking the simple pole. It gives
\[
 \int C_m C_n\,d\mu_c
 =\tfrac12\bigl(\sec(\alpha(m+n))+\sec(\alpha(m-n))\bigr).
\]
Multiplying by $\cos(\alpha m)\cos(\alpha n)$ and combining the two fractions gives the identity.
\end{proof}

\begin{proof}[Proof of Lemma~\ref{lemma:characteristic-products}]
We reduce the two determinant densities to symmetric pairs before evaluating their characteristic products.
\paragraph{Reduction by interlacing.}
It suffices to estimate with $g$ replaced by any $C_t$, $0<t<L$, uniformly in $t$. The row and column spaces for the point density are then respectively $(F_0,G_t)$ or $(A,G_t)$, where $A=\{f\in F:f(0)=0\}$ (the $-e_0 c_3$ factor extends the first determinant at column 0). They contain symmetric subspace pairs $(M,M)$ or $(U,U)$ respectively, $U=\{f\in M:f(0)=0\}$.

Here are details justifying comparison to the determinant-product densities of those symmetric pairs. For spaces $V,B$ of common dimension as above, consider the roots of
\[
 \det \int V B^T(x-\lambda)d\mu'(x),\qquad d\mu'=d\mu/(x+z^2)
\]
(function bases as column vectors in the Gram formula). The expectation in \eqref{eq:legacy-P-8} is $\prod_j r(\lambda_j)$. The roots for $(F_0,G_t)$ interlace those for $(M,M)$; similarly we can successively interlace along the chain
\[
 (U,U)\ \,\subset\ \, (F_0\cap\{f(0)=0\},M)\ \,\subset\ \, (A,G_t).
\]

To see this, each space along the chains is Chebyshev on the positive interval counting multiplicities, with a nested Chebyshev flag of all lower dimensions there (truncate jets, retaining constant first, then impose the zero at 0 when present). Gram determinants are nonzero by determinant integration. A nonzero real analytic function orthogonal to such a dimension-$m$ space must have at least $m$ sign changes: otherwise use the flag to put simple zeros of a test function just at its changes of sign, with no other zeros. Thus a real root $\lambda$ is positive: for such a root there is a nonzero real $a\in V$ with $(x-\lambda)a\perp B$, whereas $a$ by itself has at most $m-1$ flips ($m=\dim V=\dim B$). Non-real roots cannot occur either. Otherwise for some nonzero complexified $a\in V$, $(x-\lambda)a$ is orthogonal to $B$. For almost every phase $\theta$, its rotated real part needs at least $m$ sign changes. But averaging on any finite grid of positive points off zeros of $a$, the mean number of flips of $\operatorname{Re}(e^{i\theta}(x-\lambda)a)$ is bounded by the mean for $\operatorname{Re}(e^{i\theta}a)$ plus $\pi^{-1}$ times the variation of $\arg(x-\lambda)$ on the whole positive ray (a pair contributes its circular angle distance divided by $\pi$). This is bounded by $m-1$ plus a number strictly less than 1. Exhausting a dense grid gives a contradiction. Moreover consecutive pairs of nested spaces in the chains cannot share a positive root: take a real root function $a$ of the smaller first space and a real dual root function $h$ of the larger second space. By the same sign-change argument $h$ cannot belong to the smaller second space. But $(x-\lambda)a$ is orthogonal also to $h$, thus to the full larger second space, again impossible by signs.

Now continuously contract all positive jet labels to zero by common scaling. The spans limit to the ordinary consecutive polynomial spaces starting with 1 (or with $x$, when the vanishing at 0 was imposed). This convergence, also for Gram pairings, follows by taking Newton divided differences in $\beta$ (including at multiplicities), with exponential domination, subtracting constants where we impose vanishing. All roots are continuous as unordered multisets to that limit. At the polynomial limit interlacing is strict: the matrices give the usual positive-weight orthogonal polynomial roots for degrees equal to space dimensions, with weights $\mu'$ for the first chain and $x^2\mu',x\mu',x\mu'$ along the second. Consecutive degrees for a common weight interlace (symmetric multiplication compressions and no common roots by the sign argument); between the first and second weights of the latter chain use that the degree-$k-1$ polynomial for $x^2\mu'$ is proportional to $(p_k-\frac{p_k(0)}{p_{k-1}(0)}p_{k-1})/x$ for the $x\mu'$ polynomials $p_j$, so interlaces $p_k$. Here the compression characteristic polynomial is the orthogonal polynomial, e.g. by the sign/orthogonality relation $(x-\lambda)a\perp$ lower degrees. Throughout the deformation all roots remain real, and roots at adjacent levels cannot coincide. A root therefore cannot pass one from an adjacent level, so this interlacing persists. Since $\log r$ is monotone with bounded range in each application, products along adjacent interlacing levels differ by bounded factors, also in the growing-$z$ case. This proves the desired reduction to the symmetric pairs, uniformly.

\paragraph{Comparison with polynomial spaces.}
For the symmetric spaces $H=M$ or $U$, we now compare with ordinary polynomials in the original measure $\mu$. The product expectation is the determinant of the compression of the positive multiplication $r$ to $H$ in $L^2(\mu)$. It can be bracketed (up to constant factors) by polynomial compression determinants of orders $N$ and $CN$ for a sufficiently large fixed $C$. Indeed, if $\mathcal P_{<n}$ denotes polynomials of degree $<n$, then
\[
 \mathcal P_{<N}\cap\{f:f(0)=0\}\subset H,
\]
so the lower-dimensional comparison loses only one polynomial direction. For the comparison with the larger polynomial space, $M$ is approximable in the polynomials of degree $<CN$ with operator-norm error $\le e^{-N}$ in $L^2(\mu)$, as proved just below. Such approximation identifies $H$ almost isometrically with a subspace there, changing the determinant by a $1+o(1)$ factor since $r,r^{-1}$ are bounded uniformly. Determinants decrease under enlarging a space if $0<r\le1$, and increase if $r\ge1$, with codimension-one change bounded (e.g. by Schur complements).

For the approximation and polynomial estimates use the probability measures $\mu_c$ and the generating identity of Lemma~\ref{lemma:polynomial-generating}. Thus $A_m=\sum_{j\ge0} P_j(x)w_m^j$ near zero gives orthonormal polynomials $P_j$ of degree $j$. For $|m|<c/2$ this is the $L^2$ expansion also with derivatives (compare coefficients and norms by \eqref{eq:legacy-P-9}). Take $2L<c<b$. Elements $h\in M$ then have coefficients supported on $j<N$ plus sequences $q^j p_0(j)$ with $q=w_L<1$, polynomial $p_0$ of degree $<N$. For $\|h\|_{\mu_c}=1$, $|p_0(j)|\le q^{-j}$ at $j=N,\ldots,2N-1$. Lagrange interpolation bounds $|p_0(j)|\le q^{-2N}(C_1 j/N)^N$ at higher $j$, so beyond $D_0=C_2N$ the coefficients decay as $e^{-c_1 j}$. Polynomials $P_j(p^2)$ themselves are bounded by $C_3^{j+1} e^{cp/4}$ by taking a sufficiently small circle in the generating formula for $w_m$. Hence the polynomial projection through order $D_0$ has negligible $\mu_c$-norm tail at $p>C_4N$ for $C_4$ large; $h$ itself has a positive fraction of norm on $p\le C_4N$. There $\mu\gtrsim e^{-C_5 N}\mu_c$, while globally $\mu\lesssim\mu_c$. Projecting through degree $CN$ now with still larger $C$ gives the asserted exponentially good relative error in $L^2(\mu)$.

\paragraph{Evaluating the polynomial determinants.}
It remains to evaluate polynomial compression determinants. For polynomial space of degree $<n$, denote the product ratio/expectation by $R_{W,n}(r)$ for base weight $W$. One has
\[
 R_{W,n}(r^2)=K_{n,W'}(-y^2)/K_{n,W'}(-z^2),\quad
 K_{n,W'}(x)=\sup_{\deg h\le n}\frac{|h(x)|^2}{\int |h|^2dW'},\quad dW'=dW/(x+z^2)^2
\]
(using weights as measures). This is the Gram cofactor identity on degrees $\le n$ for the hyperplanes spanned by $(x+y^2)x^j$ and $(x+z^2)x^j$, $0\le j<n$, whose coefficient minors are evaluation vectors up to sign. For $c=b$ we have $K_{n,\mu_c}(-y^2)\asymp n^{2c y/\pi}$ when $y>0$ is fixed. Indeed sum $|P_j(-y^2)|^2$ via their even generating series in $u$ with $w_m=-u^2$,
\[
 (1-u^2)^{-1/2}\cosh(2s\operatorname{arctanh}u),\qquad s=cy/\pi.
\]
On the unit disk near $u=\pm1$ its absolute value is comparable to $|1\mp u|^{-s-1/2}$, elsewhere near the circle bounded. Thus the squared-coefficient sum with radial weights $r_0^{4j}$ grows as $(1-r_0)^{-2s}$ by circle integration. For truncation take $r_0=\exp(-1/n)$ (upper), then $r_0=\exp(-A/n)$ for a large fixed $A$ (lower, tail beyond $n$ bounded by $C n^{2s}\sum_{k\ge0}2^{2ks} e^{-4\cdot2^k A}\ll (n/A)^{2s}$ by the upper estimate).

For fixed positive $y,z$ and $W=\mu$ or $r\mu$, $W'$ is bounded between constants times $\mu_b/(x+Y^2)^2$ and $\mu_b$, for a fixed $Y>\max(y,z)$. We still have $K_{n,W'}(-v^2)\asymp n^{2s_v}$ for $v=y,z$, where $s_v=bv/\pi$. For the needed upper bound divide the polynomials of degree $\le n$ by $x+Y^2$ and use the $\mu_b$ norm on the resulting span (degrees $<n$ plus $g_Y=1/(x+Y^2)$). The coefficients $b_j$ of $g_Y$ on the orthonormal $P_j$ above ($c=b$) are
\[
 b_j=\frac1Y\int_0^\infty e^{-Ym}\frac{(-\tanh^2(\alpha m))^j}{\cosh(\alpha m)}dm,\qquad
 |b_j|\asymp (j+1)^{-1/2-s_Y}.
\]
Indeed $g_Y(x)=Y^{-1}\int_0^\infty e^{-Ym}\cos(m\sqrt x)dm$, and \eqref{eq:legacy-P-9} extends to imaginary $m$ paired against parameters near zero. The bounds follow substituting $v'=e^{-2\alpha m}$ (integral comparable to that of $v'^{s_Y-1/2}((1-v')/(1+v'))^{2j}$ over $(0,1)$, split near $v'\asymp1/(j+1)$ or use exponential decay above this scale). At $-v^2$, dyadic Cauchy--Schwarz and the polynomial kernel bound give $\sum_{j\ge n}|b_j P_j(-v^2)|=O(n^{s_v-s_Y})$. The full sum evaluates to $g_Y(-v^2)$ by inserting the integral, using this absolute bound (integral signs fixed for each coefficient) and the generating series at $u=\tanh(\alpha m)$. Thus after orthogonal projection onto degrees $<n$ the residual of $g_Y$ has evaluation $O(n^{s_v-s_Y})$ and its norm is at least $(\sum_{j\ge n}|b_j|^2)^{1/2}\gtrsim n^{-s_Y}$. This bounds the extra direction as required.

Consequently for $n=N$ or of order $CN$, $R_{W,n}(r^2)\asymp N^{2b(y-z)/\pi}$. Cauchy--Schwarz gives the upper bound for $R_{\mu,n}(r)$ with half exponent, and the matching lower since $R_{\mu,n}(r^2)=R_{\mu,n}(r)R_{r\mu,n}(r)$. When $y^2=2z^2=2N^{3/2}$, $W=\mu$, $W'\gtrsim z^{-4}\mu_{b+\varepsilon}$ for fixed $\varepsilon>0$, and $K_{n,W'}(-z^2)\gtrsim z^4$ by testing with a constant. The analogous generating series ($c=b+\varepsilon$) on $|u|=\exp(-1/n)$ gives $K_{n,\mu_{b+\varepsilon}}(-y^2)\le\exp(O((y+1)\log n))$, so $R_{\mu,n}(r)\le\exp(o(N))$. This proves all the product claims.
\end{proof}

In particular under either $\eta$ the probability of having $x_{N-1}<N^{3/2}$ is exponentially small (the growing-parameter product there is $\ge (3/2)^{N-1}$).

\subsection{Interpolation residuals}

The characteristic products will control the errors at $0$ and $-4$ once the following pointwise comparison is established.

\begin{auxlemma}[Residual product bounds] Take $e_{-a}$ for $a=0,4$, or $E_{-a}$ for $a=4$, and denote the relevant node set ($X$ or $\{0,X\}$) in increasing order by $v_i$. Write $\lambda_*=\pi/(3L)=8/3$. For large $N$ the residual is nonnegative on either $h_j$ and bounded above by
\begin{equation}
 C\prod_{i\text{ among first }N}\frac{v_i+a}{v_i+\lambda_*^2}.
                                                                    \label{eq:legacy-P-10}
\end{equation}
On $h_0$ there is also the lower bound with positive constant times the analogous product over all nodes. The constants are uniform in the ordered interpolation nodes.
\end{auxlemma}
\begin{proof}
First we reduce jet interpolation to monomial interpolation. To do so, translate to $w=x+a$ with evaluation at 0 (all interpolation nodes strictly positive after translation). In this coordinate each residual against the jet space is a convex combination of interpolation errors against monomial systems whose powers include $0,1,\ldots,N-1$, of the same size. Indeed take those powers first as basis polynomials; then in the Taylor expansion of the other basis functions the minors needed for Cauchy--Binet correspond to powers $J=(N\le j_1<j_2<\cdots)$ of size $N+1$ or $N+2$, using the nonzero labels $\beta=((L/2)^2,L^2,\ldots,L^2)$ and possibly $(3L/2)^2$. Write $D_J(x)$ for the determinant of the corresponding Taylor coefficients at $x$ (these jet rows evaluated by derivative orders $J$ in $x$ divided by factorials). All $D_J(0)>0$ by power positivity. Also $D'_J=\sum_r(j_r+1)D_{J+\mathbf 1_r}$, omitting collisions, and
\[
 \frac{D'_J(0)}{D_J(0)}\le\frac{\sum \beta_i}{2(2N+1)}.
\]
Indeed at 0 the columns are jets of $\beta^{j_r}/(2j_r)!$; for pure powers without factorials the sum of incremented determinants is exactly $\sum\beta_i$ times the old determinant (first at distinct labels by expansion, then take the confluent limit). Writing $\theta_N=\sum\beta_i/[2(2N+1)]$, the same estimate recursively bounds $D_J^{(k)}(0)$ by $\theta_N^kD_J(0)$. Hence
\[
 |D_J(-a)-D_J(0)|\le(e^{a\theta_N}-1)D_J(0)<D_J(0)
\]
for all large $N$, since $a\theta_N<\log2$. Thus the translated minors remain positive. Cauchy--Binet in the evaluation and extended error determinants, absolutely convergent, now gives the convex combination (weights proportional to these positive minors times ordered power determinants).

The translated $h_j(w-a)$ are positive constants times Laplace transforms of nonnegative random variables. In fact $D(x)$ has product with simple negative zeros $-s_k$, where $s_k$ range over squares $((2\pi n+\pi/3)/L)^2$, $n\in\mathbb Z$, the smallest $s_1=\lambda_*^2$. The numerator $c_1$ has negative zeros $-b_k$ with $b_k=((2n+1)\pi/L)^2$, $n\ge0$. Both products have no other factors after normalization (e.g. sine/cosine product after factoring $2\cosh t-1$ by $\sinh((t\pm i\pi/3)/2)$). Thus for $h_0$ the variable is a sum of independent exponentials of rates $s_k-a$. For $h_1$ pair each numerator factor with a distinct smaller $s_k$ (e.g. $n\ge0$ in its parametrization), so the corresponding summands are mixtures of the same exponential and zero; infinitely many other factors are unpaired. These sums converge and their survival probabilities are bounded by $C e^{-(\lambda_*^2-a)t}$, by splitting off the smallest rate (the remaining sum has finite exponential moment there). For $h_0$ the survival is also bounded below by $e^{-(\lambda_*^2-a)t}$.

For a completely monotone function in $w$ of this type, the error at 0 in a monomial system with nonnegative integer powers increases when we increase powers coordinatewise while keeping them ordered. Here is the determinant argument. All ordered power determinants with a Laplace transform row prepended are nonnegative: check for $e^{-tw}$ by multiplying columns by $e^{tw}$, expanding Taylor minors with first row the constant. The coefficient matrix of $e^{tw}w^j$ is totally nonnegative (exponential of the upper shift times $t$, a limit of products of nonnegative bidiagonals). Thus the coefficient at position $i$ (starting with 0) in the old interpolant has sign $(-1)^i$, allowing zero. If only that power is increased keeping its rank, the change in error is its coefficient times the new residual on the old power, whose sign is likewise $(-1)^i$ by power determinants including at node 0. This proves monotonicity. It brackets the error between that for ordinary consecutive powers of full size and for just powers $0,\ldots,N-1$ at the smallest $N$ nodes (for the upper bound send the highest powers successively to infinity, each then fitting only the largest remaining node after normalization). It remains to evaluate the consecutive-power errors. Write
\[
 h_j(w-a)=h_j(-a)\,\mathbb E[e^{-wZ_j}]
\]
for the nonnegative random variable supplied by the Laplace representation. At translated nodes $w_i=v_i+a$, let $Y$ be an independent sum of exponentials of rates $w_i$. For the input $e^{-tw}$, Lagrange interpolation and ordinary partial fractions show that the error at zero is $\mathbb P(Y\le t)$. Consequently the error on $h_j(w-a)$ is
\[
 h_j(-a)\,\mathbb P(Y\le Z_j)
 =h_j(-a)\,\mathbb E[\mathbb P(Z_j\ge Y\mid Y)].
\]
The survival bounds for $Z_j$ now give
\[
 \mathbb E[e^{-(\lambda_*^2-a)Y}]
 =\prod_i\frac{v_i+a}{v_i+\lambda_*^2}.
\]
Together with the monomial comparison, this proves \eqref{eq:legacy-P-10} and the lower bound.
\end{proof}

Under the measures needed in \eqref{eq:legacy-P-7}, the omitted factors in the upper-product truncation of \eqref{eq:legacy-P-10} have $v_i\ge N^{3/2}$ on the event $x_{N-1}\ge N^{3/2}$. There are $O(N)$ such factors, so the ratio between the truncated and full products is at most $\exp(O(N^{-1/2}))$. The complementary event has exponentially small probability by Lemma~\ref{lemma:characteristic-products}, and the pointwise upper bound is uniformly bounded there. For $a=0$ we can bound above with $a$ replaced by a small positive constant in these products, then send it to zero in the exponent estimate. Applying \eqref{eq:legacy-P-8} therefore proves all the $h_i$ estimates of \eqref{eq:legacy-P-7} ($b\lambda_*/\pi=1$, $b(\lambda_*-2)/\pi=1/4$).

\subsection{The normalizing coefficient}
Finally consider $O(X)=-e_0 c_3>0$. This is increasing in the ordered nodes by jet positivity: when just one node of rank $i$ (starting at 0) is increased without passing its neighbor, the change of $e_0 c_3$ equals the old interpolation weight at 0 of that node (sign $(-1)^i$) times the new error on $c_3$ at the old node (sign $(-1)^{d+i}$). Fix any $G_t$ as before in place of $G$. The expectation of $O$ with base weight $\mu$ is bounded below by that with $\mu_b$ of \eqref{eq:legacy-P-9}. Indeed the density ratio of the base weights is increasing: up to constants it is
\[
 \frac{z(\cosh(4z)+\cosh(2z)-\cosh(3z))}{\sinh(4z)},\qquad z=L\sqrt{x},
\]
with increasing numerator/denominator series coefficient ratios $\frac{2j+1}{4}(1+(1/4)^j-(9/16)^j)$.

\begin{lemma}[Increasing weight comparison]\label{lemma:weight-association}
For an ordered determinant-product probability density formed from the complete positive jet spaces used above, multiplying the one-particle measure by a positive increasing function increases the expectation of every bounded coordinatewise increasing function. The same holds for nonnegative functions by monotone truncation.\end{lemma}
\begin{proof}
The comparison is an ordered-density form of the continuous positive-association principle; see \cite[Theorem~2.27]{Bogso2014} for the general MTP$_2$ formulation and its attribution to Sarkar. We prove the required monotone coupling directly on the ordered support.

For clarity, this weight comparison uses positive association in the ordered determinant-product density, justified as follows. Discretize on a finite grid in a compact positive interval with strictly ordered tuples. In each coordinate the conditional law moves upwards stochastically if the other coordinates are raised and the weight tilted by an increasing positive base factor. To see it, supports are intervals with endpoints moving upwards. For two choices $s<t'$ in the overlap, the ratio of determinant factors (choice $t'$ over choice $s$) is increasing in each other coordinate: for each determinant separately, fix all columns except those of the two changing coordinates. Write $D(u,v)$ for the determinant with these fixed columns in their original order and the two displayed columns appended. The two-by-two Plücker relation is $D(s,a)D(t',b')-D(t',a)D(s,b')=D(s,t')D(a,b')$ for an increase $a\to b'$ in the other slot. Restoring ordered columns changes the same sign in each factor on a given side, so the right side is positive whenever all trial tuples are ordered. Raising coordinates one at a time, starting with the rightmost one, keeps these intermediate tuples ordered. The tilt only reinforces this ratio ordering. Thus single-coordinate heat-bath updates can be coupled monotonically between the distributions by common quantiles. The random-coordinate updates on each strict-tuple grid converge to equilibrium by finite irreducibility and positive holding, proving the expectation comparison. Compact grid approximation (confluent extension at collisions, or using bounded truncations) and exhaustion give the continuous comparison.
\end{proof}

Under $\mu_b$ the mean interpolant at 0 on $A_m$ of \eqref{eq:legacy-P-9} is rational in $w=w_m$ with denominator $Q_B(w)=\prod(1-w b_i')$, numerator degree $<d$ ($b_i'$ are $w$-labels of $G_t$ with multiplicities). This follows by preserving the pairings against $G_t$ and \eqref{eq:legacy-P-9}, taking jets. At the labels $a_i'$ of $F_0$ it matches $A_m(0)=1/\sqrt{1+w}$ with multiplicities. These facts still compute the mean residual at $w=1$, $m=3L/2$ by cross-integrability. Use the mixture of $1/(s+w)$, $s>1$, with density $1/(\pi\sqrt{s-1})$ for $1/\sqrt{1+w}$. Rational interpolation gives the absolute residual for each integrand at $w=1$ (all with sign $(-1)^d$)
\[
 \frac{1}{1+s}\prod_i\frac{1+s b_i'}{1-b_i'}\prod_i\frac{1-a_i'}{s+a_i'}.
\]
Since all labels are in $[0,1/3]$, identical between the multisets except for one, this is bounded below uniformly on $1<s<1+1/d$. Thus the mean $O(X)\gtrsim d^{-1/2}$, transferring to $\mu$ by the comparison and then to $G$ by positive mixture. This completes \eqref{eq:legacy-P-7}. Substitution into the finite functional of Theorem~\ref{thm:strip-functional} gives $B_N\asymp N^{-1/4}$, the strip-mass assertion of Theorem~\ref{thm:finite-bridges}.

